\subsection{投射模、内射模与扩张模}

命题10. —— 设 M 是一个 A-模。以下条件等价：

\begin{enumerate}
\def\labelenumi{(\roman{enumi})}
\item
  M 是投射的。
\item
  \(\operatorname{Ext}_{\mathbf{A}}^{i}(\mathbf{M},\mathbf{N}) = 0\) （对于每个 A-模 \(\mathbf{N}\) 以及每个整数 \(i > 0\)）.
\item
  \(\mathrm{Ext}_{\mathbf{A}}^{1}(\mathbf{M},\mathbf{N}) = 0\) （对于每个 A-模 \(\mathbf{N}\)）.
\item
  存在一个正合序列
\end{enumerate}

\[
0 \rightarrow \mathrm{K} \xrightarrow {u} \mathrm{P} \xrightarrow {r} \mathrm{M} \rightarrow 0,
\]

\(\text{其中 P 是投射的，且其中 } \text{ Ext}_{\mathbf{A}}^{1}(\mathbf{M}, \mathbf{K}) = 0.\)

\begin{enumerate}
\def\labelenumi{(\roman{enumi})}
\item
  \(\Rightarrow\) (ii)：这是 X 中命题 5 的推论，第 88 页。
\item
  \(\Rightarrow\) （iii）：这是平凡的。
\item
  \(\Rightarrow\) (iv)：这是显然的，因为 M 是自由模 P 的商模。
\item
  \(\Rightarrow\) (i)：由于 \(\operatorname{Ext}_{\mathbf{A}}^{1}(\mathbf{M},\mathbf{K}) = 0\) ，典范映射
\end{enumerate}

\[
\operatorname{Hom} _ {\mathbf {A}} (\mathbf {M}, \mathbf {P}) \rightarrow \operatorname{Hom} _ {\mathbf {A}} (\mathbf {M}, \mathbf {M})
\]

是满射（X，第 90 页，推论）；因此存在 v 的一个 A-线性截面，且 M 同构于 P 的一个直和项，从而是投射的。

命题11. --- 设 N 为一个 A-模。以下条件等价：

\begin{enumerate}
\def\labelenumi{(\roman{enumi})}
\item
  N 是内射的。
\item
  \(\operatorname{Ext}_{\mathbf{A}}^{i}(\mathbf{M},\mathbf{N}) = 0\) （对于每个 A-模 M 和每个整数 \(i > 0\)）.
\item
  \(\operatorname{Ext}_{\mathbf{A}}^{1}(\mathbf{M},\mathbf{N}) = 0\) （对于每个 A-模 \(\mathbf{M}\)）.
\item
  存在一个正合序列
\end{enumerate}

\[
0 \rightarrow \mathrm{N} \xrightarrow {u} \mathrm{I} \xrightarrow {v} \mathrm{C} \rightarrow 0,
\]

\(\text{其中 I 是内射的且其中 } \text{ Ext}_{\mathbf{A}}^{1} (\mathbf{C}, \mathbf{N}) = 0;\)

\begin{enumerate}
\def\labelenumi{(\alph{enumi})}
\setcounter{enumi}{21}
\item
  \(\operatorname{Ext}_{\mathbf{A}}^{1}(\mathbf{M},\mathbf{N}) = 0\) （对于每个单演 A-模 \(\mathbf{M}\)）.
\item
  \(\Rightarrow\) (ii)：这是 X 中命题 5 的推论，第 88 页。
\end{enumerate}

\begin{enumerate}
\def\labelenumi{(\roman{enumi})}
\setcounter{enumi}{1}
\item
  \(\Rightarrow\) (iii) \(\Rightarrow\) (v)：这是平凡的。
\item
  \(\Rightarrow\) (iv)：这是显然的，因为 N 是内射模的子模（X，第 19 页，推论 3）。
\item
  \(\Rightarrow\) (i)：由于 \(\operatorname{Ext}_{\mathbf{A}}^{1}(\mathbf{C},\mathbf{N}) = 0\) ，典范同态
\end{enumerate}

\[
\operatorname{Hom} _ {\mathbf {A}} (\mathrm{I}, \mathbf {N}) \rightarrow \operatorname{Hom} _ {\mathbf {A}} (\mathbf {N}, \mathbf {N})
\]

是满射的（X，第 93 页，推论）；因此存在 u 的一个 A-线性收缩，且 N 同构于 I 的一个直和项，从而 N 是内射的（X，第 16 页，命题 9）。

\begin{enumerate}
\def\labelenumi{(\alph{enumi})}
\setcounter{enumi}{21}
\tightlist
\item
  \(\Rightarrow\) (i)：若 a 是 A 的一个理想，则有 \(\operatorname{Ext}_{\mathbf{A}}^{1}(\mathbf{A}/\mathfrak{a},\mathbf{N})=0\) ；因此典范映射 \(\operatorname{Hom}_{\mathbf{A}}(\mathbf{A},\mathbf{N})\to\operatorname{Hom}_{\mathbf{A}}(\mathfrak{a},\mathbf{N})\) 是满射的，且 N 是内射的（X，第 16 页，命题 10）。
\end{enumerate}

